% Introduction and Related Work for the writing agent to integrate.
% Unified narrative voice: declarative, quantitative, cautious about causal claims.

\section{Introduction}

Human beliefs shape behavior and wellbeing, yet how a subjective mental construct is mapped onto neurobiological activity in a quantitatively graded fashion is largely unresolved. This gap is especially consequential for substance use disorders, where purely pharmacological accounts cannot explain the heterogeneity of drug response and the durability of craving after abstinence~\cite{picciotto2008nicotinic,benowitz2009pharmacology}. The placebo effect is the canonical case for belief-to-brain interaction: positive expectations about treatment can alter clinically meaningful outcomes and produce measurable shifts in neural activity even in the absence of active pharmacology~\cite{wager2004placebo,benedetti2013placebo}. Within addiction neuroscience, work with nicotine-dependent smokers has shown that binary beliefs about whether a cigarette contains nicotine selectively modulate reward prediction error signals in the ventral striatum~\cite{gu2015belief} and alter insular responses in deprived smokers~\cite{gu2016belief}. However, these demonstrations have been restricted to an all-or-none operationalization of belief, leaving open the question of whether beliefs can modulate neural activity with the graded precision that defines pharmacological dose-response relationships.

Pharmacology quantifies the impact of an agent through the dose-response function, in which the amount of active ingredient in a drug scales proportionally to a downstream biological readout. Nicotine in particular exerts well-characterized dose-dependent effects on regional brain activity measured with PET and functional MRI~\cite{brody2006cigarette,cosgrove2009beta2}. The thalamus is an especially informative locus in this regard: it harbors one of the highest densities of nicotinic acetylcholine receptors (nAChRs) in the human brain~\cite{picciotto2008nicotinic,cosgrove2009beta2}, and its activity has been shown to scale with nicotine exposure~\cite{brody2006cigarette}. At the same time, PET imaging has shown that smoking denicotinized or low-nicotine-content cigarettes is sufficient to occupy a substantial fraction of nAChRs in the human brain~\cite{brody2011ventral,cosgrove2009beta2}, implying that factors beyond raw nicotine dose---plausibly including expectations, habits, and beliefs about the act of smoking---shape receptor-level readouts.

We reasoned that if subjective beliefs about a drug can act on the brain with quantitative precision, then an experimental manipulation that holds actual nicotine dose constant while parametrically varying the instructed belief about the dose should yield graded, dose-like neural responses in thalamic circuitry. To test this, we instructed 20 nicotine-dependent human smokers (60 fMRI scans in total) across three visits that an e-cigarette they were about to vape contained ``low'', ``medium'', or ``high'' levels of nicotine, while actual nicotine content was held constant at 1.2\% across all cartridges. After vaping, participants performed a value-based investment task during fMRI. A non-smoking healthy-control group (n=31) was scanned on the same task without vaping and was used as a reference. We chose e-cigarettes as the delivery vehicle because they permit finer control over nicotine strength than traditional cigarettes and because smokers in our sample were naive to vaping, reducing conditioned cue effects that have been proposed to drive striatal responses in prior belief-by-drug experiments~\cite{gu2015belief}.

Our predictions were grounded in the anatomical and pharmacological properties of thalamocortical circuits. Because the thalamus gates sensory and value-related information and is densely innervated by nAChRs~\cite{sherman2016thalamus,halassa2017thalamic}, we hypothesized that value-related thalamic activity would scale parametrically with instructed belief dose. We further predicted that thalamocortical coupling with the ventromedial prefrontal cortex (vmPFC)---a region repeatedly linked to value representation~\cite{plassmann2007orbitofrontal,levy2012root}, task-state encoding~\cite{wilson2014orbitofrontal,schuck2016human}, and belief-related inference~\cite{rouault2018belief,rouault2022prefrontal}---would likewise track belief dosage. In contrast, we expected that the ventral striatum, which is strongly shaped by conditioned cigarette cues and reinforcement history~\cite{haber2010reward,schultz1997neural,kober2010brain}, would track reward value overall but would not differentiate between belief conditions in a cohort unfamiliar with the delivery device.

\section{Related Work}

\paragraph{Beliefs and the brain in addiction.} Prior work established that telling smokers a cigarette contained or did not contain nicotine selectively shifted striatal value-learning signals despite identical pharmacology~\cite{gu2015belief}. A follow-up study showed that the same manipulation altered insular activity and subjective craving in abstinent smokers~\cite{gu2016belief}. These studies framed belief as a binary factor (nicotine present vs. absent). Our design extends this line of work by introducing a three-level dose manipulation and by probing thalamic rather than only striatal circuitry.

\paragraph{Dose-response imaging of nicotine.} Pharmacological imaging has established dose-dependent thalamic and striatal effects of nicotine itself~\cite{brody2006cigarette,brody2011ventral,cosgrove2009beta2}. Across these studies, the thalamus has emerged as the region most reliably tracking nicotine exposure, consistent with its high nAChR density~\cite{picciotto2008nicotinic}. Our study asks the complementary question: whether a purely cognitive manipulation can produce an analogous graded effect in the same anatomical substrate.

\paragraph{Placebo and expectation.} Expectation-based placebo effects recruit a distributed circuitry including prefrontal, striatal, and brainstem regions~\cite{wager2004placebo,benedetti2013placebo}. Most placebo experiments contrast an active and an inactive condition, leaving unclear whether expectation acts on the brain with pharmacological-style precision. We bring the dose-response logic of pharmacology into the placebo tradition by using three graded levels of instructed strength.

\paragraph{Value, state, and belief in vmPFC.} The vmPFC encodes subjective value across reward modalities~\cite{plassmann2007orbitofrontal,levy2012root} and represents abstract task states~\cite{wilson2014orbitofrontal,schuck2016human}. Recent work has shown that the vmPFC participates in self-performance estimation and belief updating~\cite{rouault2018belief,rouault2022prefrontal}. These findings motivate our focus on thalamocortical coupling with vmPFC as the circuit most likely to carry parametric belief signals.

\paragraph{Thalamus in cognition.} The thalamus is now appreciated as more than a passive relay: it contributes actively to cortical computation and cognitive control~\cite{sherman2016thalamus,halassa2017thalamic}. Its role in gating incoming information makes it a plausible substrate for belief-driven dose-like modulation, particularly given its nAChR density~\cite{picciotto2008nicotinic}.

\paragraph{Methodological substrate.} For atlas-based ROI definition we used the WFU Pick Atlas~\cite{maldjian2003automated} and the THOMAS thalamic parcellation~\cite{mcleod2017spatially}. Motion artefact was mitigated with ArtRepair~\cite{mazaika2009artifact}, and connectivity was evaluated with psychophysiological interaction analysis~\cite{friston1997psychophysiological}. Participant characterization used the Fagerstr\"{o}m Test for Nicotine Dependence~\cite{fagerstrom1991measuring} and companion substance-use questionnaires.
